% ====================================================================
% PROPOSED MANUSCRIPT TEXT: TASK 6 OOD GENERALIZATION (DO NOT APPLY DIRECTLY)
% Prepared strictly from verified evidence in review_package_20261023_172944Z
% ====================================================================

\subsection{Controlled Out-of-Distribution Phase-Shift Generalization}
To evaluate kernel generalization under distribution shift across a macroscopic quantum phase boundary, we partition ground states of the 1D disordered transverse-field Ising model ($n = 8$ spins, open boundary conditions) into non-overlapping parameter regimes. Models are trained on ground states generated in the higher-field paramagnetic and critical ensemble ($h_i \sim \mathcal{U}(0.8, 1.5)$, $N_{\text{train}} = 40$) and evaluated on a held-out lower-field ensemble deep within the ordered ferromagnetic regime ($h_i \sim \mathcal{U}(0.2, 0.7)$, $N_{\text{test}} = 100$). The learning target is half-chain bipartite entanglement entropy $S_{n/2}(\psi)$. This benchmark is conducted as a separate noiseless evaluation distinct from the finite-shot budget trade-off experiments.

In the out-of-distribution regime, the small total variance of entanglement entropy across ferromagnetic ground states ($\operatorname{Var}(y_{\text{ood}}) \approx 4 \times 10^{-4}$) increases metric sensitivity to residual biases: small absolute offsets produce negative out-of-distribution $R^2$ values across all evaluated architectures (mean $R^2_{\text{ood}} = -3.26$ for FQ vs $-70.4$ for PQK-2nn). However, evaluation in absolute physical units reveals that the full quantum fidelity kernel (FQ) achieves substantially lower mean absolute error than local correlator architectures: FQ achieves an out-of-distribution MAE of $0.0602 \pm 0.0085$ compared to $0.1850 \pm 0.0294$ for PQK-2nn (mean paired difference $-0.1248 \pm 0.0334$; two-sided Wilcoxon signed-rank test $W = 21.0$, $p = 8.51 \times 10^{-4}$, Holm-adjusted $p = 8.51 \times 10^{-4}$ across $N = 20$ independent disorder realizations). At the individual disorder level, FQ achieves lower MAE in 16 of the 20 seeds, while PQK-2nn achieves lower MAE in 4 seeds (Seeds 8, 9, 11, and 16). FQ also significantly outperforms constant training-set predictors (training-mean baseline MAE $0.6061 \pm 0.0075$ and training-median baseline MAE $0.6121 \pm 0.0086$, both Holm-adjusted $p = 5.72 \times 10^{-6}$).

A separate within-model decomposition of the fitted support vector regressor (SVR) explains this performance. Because the training ensemble lies in the paramagnetic regime where entropy is moderate ($y \approx 0.392$), a constant intercept-only predictor fails out-of-distribution, incurring an MAE of $0.4616 \pm 0.0162$. The signed kernel contribution $g(x) = \sigma_{\text{train}} \sum_{i} \alpha_i K(x_i, x)$ provides a regularized positive offset ($+0.4379$), shifting predictions into the ferromagnetic target domain ($y \approx 0.993$) and improving MAE in all 20 evaluated seeds (mean improvement $-0.4014 \pm 0.0156$). This decomposition represents an algebraic property of SVR regularization in reproducing kernel Hilbert spaces with decaying cross-domain fidelity overlaps ($K \approx 0.93$ in-distribution to $K \approx 0.43$ out-of-distribution), which prevents extrapolative runaway without asserting a causal physical discovery mechanism.
